% --- INICIO DE PLANTILLA PERSONALIZADA ---
\documentclass[oneside,11pt]{Latex/Classes/PhDthesisPSnPDF}

% --- Paquetes Originales ---
\usepackage{blindtext}
\usepackage{actuarialangle} 
\usepackage{tabularx}
\usepackage{float}
\usepackage{multirow, array}
\usepackage[usenames,dvipsnames,svgnames,table]{xcolor}
\usepackage{longtable}
\usepackage{latexsym,amsmath,amssymb,amsfonts}
\usepackage{enumitem}
\usepackage{xparse}
\usepackage{booktabs}
\usepackage[round, sort, numbers]{natbib}  
\usepackage{adjustbox}
\usepackage[utf8]{inputenc}
\usepackage{caption}
\usepackage{graphicx}
\usepackage{hyperref}
\usepackage{cleveref}

% Definición de colores
\definecolor{miblue}{RGB}{68, 114, 196}

% --- CAMBIO CRÍTICO 1: Usar \input para macros en el preámbulo ---
% \include da problemas antes del begin document. \input es seguro.
\input{Latex/Macros/MacroFile1}

% --- Definiciones de seguridad (por si la clase falla al cargar alguna variable) ---
\providecommand{\facultad}[1]{\def\lafacultad{#1}}
\providecommand{\escudofacultad}[1]{\def\elescudo{#1}}
\providecommand{\degree}[1]{\def\elgrado{#1}}
\providecommand{\director}[1]{\def\eldirector{#1}}
\providecommand{\lugar}[1]{\def\ellugar{#1}}
\providecommand{\degreedate}[1]{\def\lafecha{#1}}

% --- Configuración para bloques de código de R (Pandoc) ---

%%%%%%%%%%%%%%%%%%%%%%%%%%%%%%%%%%%%%%%%%%%%%%%%%%%%%%%%%%%%%%%%%%%%%%%%%%%%%%%%
%                         DATOS (Conectados a RMarkdown)                       %
%%%%%%%%%%%%%%%%%%%%%%%%%%%%%%%%%%%%%%%%%%%%%%%%%%%%%%%%%%%%%%%%%%%%%%%%%%%%%%%%
\title{Análisis Descriptivo de los Factores Determinantes en ciclos
agrícolas}
\author{Joshua Santeliz \\ Emily Gallardo}

% Lógica para Facultad: Si la defines en el YAML la usa, si no, usa la default

\facultad{Facultad de Ciencias Económicas y Sociales \\ Escuela de
Estadistica y Ciencias Actuariales}
  \escudofacultad{Latex/Classes/Escudos/faces}

\degree{Computación I}
\director{Jesus Ochoa \\ Oliver Triveño}
\degreedate{Febrero 2026} 
\lugar{Caracas}

\portadatrue 

% Metadatos PDF
\keywords{}
\subject{}

% --- CORRECCIÓN PARA PANDOC NUEVO (Imágenes) ---
\makeatletter
\providecommand{\pandocbounded}[1]{#1}
\makeatother


%%%%%%%%%%%%%%%%%%%%%%%%%%%%%%%%%%%%%%%%%%%%%%%%%%%%%
%                   DOCUMENTO                       %
%%%%%%%%%%%%%%%%%%%%%%%%%%%%%%%%%%%%%%%%%%%%%%%%%%%%%
\begin{document}

\maketitle

%%%%%%%%%%%%%%%%%%%%%%%%%%%%%%%%%%%%%%%%%%%%%%%%%%%%%
%                  PRÓLOGO                          %
%%%%%%%%%%%%%%%%%%%%%%%%%%%%%%%%%%%%%%%%%%%%%%%%%%%%%
\frontmatter

% Puedes agregar dedicatorias desde el YAML si quieres, o dejarlas fijas aquí

%%%%%%%%%%%%%%%%%%%%%%%%%%%%%%%%%%%%%%%%%%%%%%%%%%%%%
%                   ÍNDICES                         %
%%%%%%%%%%%%%%%%%%%%%%%%%%%%%%%%%%%%%%%%%%%%%%%%%%%%%
\setcounter{secnumdepth}{3}
\setcounter{tocdepth}{3}

\tableofcontents



\cleardoublepage

  \cleardoublepage       % Empieza en página derecha (estándar en tesis)
  \chapter*{Resumen}     % Crea el título "Resumen" sin numeración (Capítulo *)
  \addcontentsline{toc}{chapter}{Resumen} % (Opcional) Lo añade al índice
  
  En el presente informe se realiza un análisis exploratorio de los
datos de influencias en los cultivos en el territorio de la India,
tomando en cuenta registros desde 2010 hasta 2020. Es un informe que
tiene como propósito responder a la incertidumbre que enfrentan los
ciclos agrícolas al ser afectados ya sea positivamente o negativamente
por factores como los agroquímicos y la precipitación anual, esto en un
determinado espacio de tierra. Se usan métodos estadísticos descriptivos
(como el uso de medidas de tendencia central, cocientes estadísticos
entre otros conceptos) con el fin de medir dicha incertidumbre por
mediante el uso de indicadores de dispersión, estos indicadores son: la
desviación típica, los rangos y el coeficiente de variación. También nos
apoyaremos en estadísticas bivariantes con el fin de comprender las
relaciones lineales entre variables, utilizando Covarianza y coeficiente
de correlación para observar su fuerza de influencia. Además, se
examinan las tendencias históricas aplicando las medidas de forma,
observando su distribución y concentración de los datos. Por último,
analizaremos el desempeño de producción por temporadas (Kharif, Zaid,
Rabi), aplicando técnicas estadísticas descriptivas y visualización de
datos con
R.             % Aquí se pega el texto que escribiste en el YAML
  
  \cleardoublepage

%%%%%%%%%%%%%%%%%%%%%%%%%%%%%%%%%%%%%%%%%%%%%%%%%%%%%
%                   CONTENIDO (El Corazón)          %
%%%%%%%%%%%%%%%%%%%%%%%%%%%%%%%%%%%%%%%%%%%%%%%%%%%%%
\mainmatter
\def\baselinestretch{1.5}

% --- CAMBIO CRÍTICO 2: Aquí RMarkdown inyecta todo tu texto ---
\chapter*{Introducción}

La industria agrícola ha experimentado a lo largo del tiempo una
dependencia importante a las temporadas de cultivo, siendo esta una de
las cosas en el cual la industria ha tenido considerables inversiones,
por ejemplo, en la Ciencia y tecnología. La ciencia es un ámbito que
ofrece alternativas basadas en datos, una de ellas es el análisis de
datos y el uso de la estadística aplicada a la agricultura. Comunicando
una comprensión a las temporadas de cultivo a fin de aplicar inversiones
acertadas y competentes que favorezcan a las producciones. Siguiendo
esta idea entonces debemos responder a la incertidumbre de como
gestionar las temporadas de cultivos para obtener resultados óptimos. Lo
que se vuelve fundamental comprender factores determinantes que afecten
a los ciclos agrícolas. Algunos de ellos son el clima y los agroquímicos
(fertilizantes y pesticidas).

Para el profesional en Ciencias Estadísticas y Ciencia de Datos, este
sector presenta un desafío fascinante: ¿Es posible modelar la
incertidumbre en temporadas de cultivo? ¿Existen patrones que le
permitan a las industrias entender los ciclos agrícolas?, bueno en este
informe trataremos de implementarlo y descubrir esos patrones de una
manera básica pero eficaz, la cual podrá entender sin necesidad de
conocimientos o conceptos complejos en Estadística.

El presente trabajo de investigación tiene como propósito realizar un
análisis estadístico exhaustivo sobre el conjunto de datos''
Agricultural Crop Yield in Indian States Dataset'', el cual abarca más
de 19,000 registros de cultivos desde 1997 hasta 2020. A través del uso
de herramientas computacionales avanzadas en R y técnicas de
visualización de datos, se busca descomponer la varianza de los cultivos
en función de factores críticos como el Clima, los Fertilizantes y los
Pesticidas.

\section*{Estructura del Informe}

Este documento técnico se encuentra estructurado en cinco capítulos
fundamentales:

\begin{itemize}
    \item En el \textbf{Capítulo I}, se plantea la problemática de no comprender las influencias en los ciclos agrícolas, formulando las preguntas de investigación y delimitando el alcance temporal y geográfico del estudio.
    
    \item El \textbf{Capítulo II} establece el Marco Téorico, definiendo los conceptos estadísticos clave como las medidas de tendencia central, Medidas de Dispersión y correlación, necesarios para la interpretación de los datos.
    
    \item El \textbf{Capítulo III} describe el Marco Metodológico, detallando el proceso de limpieza de datos (ETL), las pruebas de hipótesis y los algoritmos utilizados para validar la calidad de la información.
    
    \item El \textbf{Capítulo IV} presenta el Análisis de Resultados, donde se exhiben las tablas y gráficos generados dinámicamente, discutiendo los hallazgos sobre la variabilidad de los ciclos agrícolas y la dinámica temporal.
    
    \item Finalmente, en el \textbf{Capítulo V}, se exponen las Conclusiones y Recomendaciones estratégicas derivadas del análisis cuantitativo, orientadas a la toma de decisiones basada en datos (\textit{Data-Driven Decision Making}).
\end{itemize}

\chapter{Planteamiento del Problema}

En el mundo agrícola hay muchos temas por tratar, desde la salud
nutricional hasta la gastronomía industrial, siendo temas de vital
importancia para la vida humana. En este estudio nos vamos a ubicar en
la agricultura de la India en el ámbito del cultivo, resaltando que una
gestión eficiente de las temporadas de cultivo no solo garantiza
cosechas óptimas, sino que genera seguridad alimentaria y eleva la
calidad de vida de las personas.

Analizaremos específicamente las temporadas de cultivo y como estas son
influenciadas por distintos factores como el clima y los agroquímicos.
Estas influencias son de gran interés para las empresas ya que una
descripción de las mismas, servirían de referencia para mejorar su
gestión de inversión.

Por lo antes mencionado surgen problemas por ejemplo a la hora de saber
que cantidad de agroquímicos aplicar a un área en específico y en cierta
estación de cultivos y también no saber el nivel de precipitación que
tendremos en un año, lo cual puede ser determinante para prepararse en
los ciclos de cultivo en dicho año.\\

Estos problemas se resuelven luego de modelar esta investigación donde
podremos comprender mejor la situación y se recomendará posibles
decisiones a tomar. Teniendo en cuenta esto surgen las siguientes
preguntas de investigación:\\

\textbf{1.} ¿De qué manera influyen la cantidad de agroquímicos en los
ciclos de cultivos?\\

\textbf{2.} ¿Por qué son necesario los fertilizantes y los pesticidas en
una temporada de cultivo?\\

\textbf{3.} ¿Cuáles son las precipitaciones anuales que favorecen a las
temporadas de cultivo?\\

\textbf{4.} ¿La cantidad de agroquímicos varía según los niveles de
precipitación?\\

\section{Justificación}

La justificación de este trabajo de investigación reside en la necesidad
de aplicar técnicas estadísticas rigurosas y herramientas de software
como R y Python para mitigar la incertidumbre de sembrar los cultivos en
el sector agrícola de la India.

Desde el punto de vista práctico, los resultados permitirán a las
Industrias tomar decisiones basadas en datos (Data-Driven Decisions)
sobre que materiales invertir o en qué temporada del año se debe
desarrollar. Con una metodología replicable utilizando R y LaTeX para el
análisis exploratorio y descriptivo de grandes volúmenes de datos.\\

\section{Objetivos}

\subsection{Objetivo General}
\begin{itemize}
  \item{Realizar un análisis estadístico descriptivo de datos de cultivos en diferentes temporadas de la India durante el período 2010-2020.  para identificar y describir las características más influyentes en las temporadas de cultivos. }
\end{itemize}

\subsection{Objetivos Específicos}
\begin{itemize}
  \item{Analizar la variabilidad pluviométrica mediante el uso de indicadores estadísticos, con el fin de determinar una óptima gestión del recurso hídrico. }
  \item{Observar la distribución de los factores que afectan a los cultivos por temporada del año. }
  \item{Categorizar temporadas de mayor y menor demanda tanto de pesticidas como de fertilizantes. }
\end{itemize}

\section{Cobertura}

\subsection{Horizontal}

Describir la cobertura Horizontal del trabajo de investigación\\

\begin{table}[ht]
\centering
\caption{Variables consideradas} 
\resizebox{10cm}{!} {
\begin{tabular}{ccc}
  \hline
\ \ \ \ \ NOMBRE DE LA VARIABLE \\ 
  \hline
   VARIABLE 1 \\ 
VARIABLE 2\\ 
VARIABLE n\\ 
   \hline
\end{tabular}
}
\end{table}

\subsection{Vertical}

Describir la cobertura Vertical del trabajo de investigación

\section{Periodo de Referencia}

Describir la fuente y el Periodo de Referencia utilizado en el trabajo
de investigación

\section{Variables}

Describir la variable en estudio en el trabajo de investigación

\chapter{Marco Teórico}

\section{Bases Teóricas}

Resumen de las bases teoricas explicadas en el capitulo y su
importancia\\

\subsection{Definición 1}

Descripción de la definición 1\\

\subsection{Definición 2}

Descripción de la definición 2\\

\subsection{Definición n}

Descripción de la definición n\\

\chapter{Marco Metódico}

En esta sección se presentan las metodologías y/o test necesarios
relacionados con la práctica divulgada en el marco teórico, así como las
fuentes de datos.

\section{Bases metódicas utilizadas en el trabajo de investigación}

Desarrollo\\

\subsection{Item 1}

Descripción\\

\subsection{Item 2}

Descripción\\

\subsubsection{Item 2.2}

Descripción\\

\subsection{Item n}

Descripción\\

\chapter{Análisis de Resultados}

\{\small Este capítulo presenta los resultados obtenidos al aplicar las
metodologías descritas en el capítulo anterior, así como del análisis
estadístico descriptivo de las variables estudiadas y las ventajas y
desventajas al aplicar dichos tests.\}

\section{Presentación de resultados 1}

Descripción\\

\section{Presentación de resultados 2}

Descripción\\

\section{Presentación de resultados n}

Descripción\\

\chapter{Conclusiones y Recomendaciones}

\textbf{1.-} Conclusion 1 del trabajo de investigación.\\

\textbf{Recomendación} Recomendación 1 del trabajo de investigación.\\

\textbf{2.-} Conclusion 2 del trabajo de investigación.\\

\textbf{Recomendación} Recomendación 2 del trabajo de investigación.\\

\textbf{n.-} Conclusion n del trabajo de investigación.\\

\textbf{Recomendación} Recomendación n del trabajo de investigación.\\

\appendix
\renewcommand{\thechapter}{\Alph{chapter}}

\chapter{Anexo A}

\chapter{Anexo B}

\chapter{Anexo C}

%%%%%%%%%%%%%%%%%%%%%%%%%%%%%%%%%%%%%%%%%%%%%%%%%%%%%
%                   REFERENCIAS                     %
%%%%%%%%%%%%%%%%%%%%%%%%%%%%%%%%%%%%%%%%%%%%%%%%%%%%%
% Mantenemos la estructura natbib de tu plantilla
\renewcommand{\bibname}{Lista de Referencias}
\bibliographystyle{apalike} 
\bibliography{bibliografia.bib} 

%%%%%%%%%%%%%%%%%%%%%%%%%%%%%%%%%%%%%%%%%%%%%%%%%%%%%
%                   APÉNDICES                       %
%%%%%%%%%%%%%%%%%%%%%%%%%%%%%%%%%%%%%%%%%%%%%%%%%%%%%

\end{document}
